\documentclass{article}
\usepackage{pxSRC}
\usepackage{pxLST}
\begin{document}
\section{Literal text}
% Literal logs are never executed as TeX.
\pxInput[as=text,first-line=2,last-line=3]{fixtures/pxSRC.recovery-text.txt}
\ifcsname pxTextMustNotRun\endcsname
  \PackageError{pxSRC-modes}{Literal input executed TeX}{}
\fi
% Registered defaults must be overridden only for this use.
\pxSource[as=text,first-line=2,last-line=3]{log}{fixtures/pxSRC.recovery-text.txt}
\pxUseSource[container=minipage,width=.75\linewidth]{log}

\section{Code}
% The original code file is rendered using pxLST, preserving its line numbers.
\pxInput[as=code,language=Python,first-line=2,last-line=4,
  title={Selected Python},container=minipage,width=.9\linewidth]
  {fixtures/pxSRC.recovery-code.py}
\ifcsname pxCodeMustNotRun\endcsname
  \PackageError{pxSRC-modes}{Source code executed TeX}{}
\fi
% Non-TeX marker styles must work without a separate extraction engine.
\pxInput[as=code,region=selected,region-prefix=//,
  title={Selected JavaScript}]{fixtures/pxSRC.recovery-marked.js}
\ifcsname pxCodeRegionMustNotRun\endcsname
  \PackageError{pxSRC-modes}{Region code executed TeX}{}
\fi

\section{Existing native heading support}
% Heading-based extraction from unmodified PR files remains supported.
\pxInput[subsection={Metodologia}]{fixtures/reuse/PR1.tex}
\ifcsname pxTestSubsection\endcsname\else
  \PackageError{pxSRC-modes}{Accepted heading selector regressed}{}
\fi
\ifcsname pxNestedInput\endcsname\else
  \PackageError{pxSRC-modes}{Source-relative native input regressed}{}
\fi
\typeout{PXSRC-MODES-PASS}
\end{document}
